\documentclass[11pt]{article}
\usepackage[margin=1in]{geometry}
\usepackage{amsmath,amssymb,bm}
\usepackage{booktabs}
\usepackage{hyperref}
\newcommand{\Rnav}{\mathcal{R}_{\mathrm{nav}}}
\newcommand{\sigmar}{\sigma_r}
\newcommand{\sigmav}{\sigma_v}
\newcommand{\GammaGeom}{\Gamma}
\newcommand{\xstate}{\bm{x}}
\newcommand{\zhat}{\hat{\bm{z}}}
\newcommand{\Phimat}{\bm{\Phi}}
\newcommand{\Qmat}{\bm{Q}}
\newcommand{\Rmat}{\bm{R}}
\newcommand{\Pmat}{\bm{P}}
\newcommand{\Kmat}{\bm{K}}
\newcommand{\Hmat}{\bm{H}}
\newcommand{\Imat}{\bm{I}}

\title{Explaining the Cislunar Navigation Risk Equation}
\author{Fabio Facin}
\date{\today}
\begin{document}
\maketitle

\section{The Idea}

The spacecraft always has an estimate of where it is, but that estimate can be
more or less trustworthy. The navigation risk equation turns several warning
signals into one score.

\begin{equation}
\Rnav(t)
=
w_r \frac{\sigmar(t)}{r_0}
+
w_v \frac{\sigmav(t)}{v_0}
+
w_g \frac{1}{\GammaGeom(t)+\epsilon}
+
w_m M(t)
+
w_l L(t).
\end{equation}

In plain language: risk goes up when the spacecraft is uncertain about
position, uncertain about velocity, has bad viewing geometry, misses sensor
measurements, or has poor lighting/visibility.

\section{Position and Velocity Uncertainty}

The EKF covariance matrix tells us how uncertain the filter is. The position
block gives
\begin{equation}
  \sigmar(t)=\sqrt{\mathrm{tr}(\Pmat_{rr}(t))},
\end{equation}
and the velocity block gives
\begin{equation}
  \sigmav(t)=\sqrt{\mathrm{tr}(\Pmat_{vv}(t))}.
\end{equation}
The scale constants $r_0$ and $v_0$ make these terms dimensionless.

\section{Geometry Strength}

The term $\GammaGeom(t)$ measures whether the sensor geometry is helpful. Good
viewing geometry makes the spacecraft easier to locate. Poor geometry makes
the same sensor noise more dangerous. Because the equation uses
$1/(\GammaGeom+\epsilon)$, weak geometry increases risk.

\section{Missed Measurements and Lighting}

$M(t)$ is a missed-measurement penalty. $L(t)$ is a lighting or visibility
penalty. These terms matter because optical navigation and ranging can degrade
during outages, glare, eclipse, poor target contrast, or sensor interruptions.

\section{How the Score Changes the EKF}

The score can adapt measurement noise:
\begin{equation}
  \Rmat(t)=\Rmat_0[1+\beta\,\mathrm{clip}(\Rnav,0,R_{\max})],
\end{equation}
or process noise:
\begin{equation}
  \Qmat(t)=\Qmat_0[1+\alpha\,\mathrm{clip}(\Rnav,0,R_{\max})].
\end{equation}
This tells the filter to become more cautious when navigation conditions are
bad. The clipping term prevents the adaptation from becoming unbounded.

\section{Normal-Person Explanation}

The formula is like a dashboard warning light for spacecraft navigation. If
the spacecraft has good sensor data and low uncertainty, the warning stays low.
If it loses measurements, has bad viewing angles, or flies through poor
lighting conditions, the warning rises. The filter can then respond by trusting
its sensors or dynamics less aggressively.

\section{Why This Is Research}

The research question is not whether the formula looks nice. The question is
whether it improves navigation in simulation compared with a normal EKF. If it
helps during outages and bad geometry without making clean cases worse, it is
useful. If it increases error or only works after arbitrary tuning, it fails.

\end{document}
